\begin{exercise}[Universal triangle]
State the universal property of a sheafification $(\mathcal F^\#,\eta)$ and explain why the factorization morphism is unique.
\end{exercise}

\begin{exercise}[Uniqueness of associated sheaf]
Prove that any two sheafifications of the same presheaf are uniquely isomorphic compatibly with their unit maps.
\end{exercise}

\begin{exercise}[Basis for the espace etale]
Prove that the sets $\widetilde s(U)=\{s_x:x\in U\}$ form a basis for a topology on $E_{\mathcal F}=\coprod_x\mathcal F_x$.
\end{exercise}

\begin{exercise}[Local homeomorphism]
Show that $\pi:E_{\mathcal F}\to X$ restricts to a homeomorphism $\widetilde s(U)\to U$ on every basic open.
\end{exercise}

\begin{exercise}[Local representability]
Prove that a section $\sigma:U\to E_{\mathcal F}$ is continuous iff it is locally of the form $x\mapsto s_x$ for sections $s$ of the original presheaf.
\end{exercise}

\begin{exercise}[Associated sections form a sheaf]
Using ordinary gluing of maps, prove directly that $U\mapsto\mathcal F^\#(U)$ is a sheaf.
\end{exercise}

\begin{exercise}[The canonical morphism]
Prove that $s\mapsto\widetilde s$ is compatible with restriction and hence defines $\eta:\mathcal F\to\mathcal F^\#$.
\end{exercise}

\begin{exercise}[Surjectivity on stalks]
Prove that $(\eta_\mathcal F)_x:\mathcal F_x\to(\mathcal F^\#)_x$ is surjective.
\end{exercise}

\begin{exercise}[Injectivity on stalks]
Prove that $(\eta_\mathcal F)_x$ is injective.
\end{exercise}

\begin{exercise}[Factorization existence]
Let $u:\mathcal F\to\mathcal G$ with $\mathcal G$ a sheaf. Construct locally and then glue the factor $\bar u:\mathcal F^\#\to\mathcal G$.
\end{exercise}

\begin{exercise}[Factorization uniqueness]
Prove the factor in the universal property is unique because every sheafified section is locally represented by original sections.
\end{exercise}

\begin{exercise}[Functoriality]
For $f:\mathcal F\to\mathcal G$, construct $f^\#$ and prove $(gf)^\#=g^\#f^\#$.
\end{exercise}

\begin{exercise}[Naturality of the unit]
Prove $f^\#\eta_\mathcal F=\eta_\mathcal Gf$.
\end{exercise}

\begin{exercise}[Stalk map of a sheafified morphism]
Under the canonical stalk identifications, prove $(f^\#)_x=f_x$.
\end{exercise}

\begin{exercise}[Sheafifying a sheaf]
If $\mathcal G$ is a sheaf, prove $\eta_\mathcal G$ is an isomorphism using only the universal property.
\end{exercise}

\begin{exercise}[Idempotence]
Deduce $(\mathcal F^\#)^\#\cong\mathcal F^\#$ canonically.
\end{exercise}

\begin{exercise}[Constant presheaf on two points]
On a two-point discrete space, compute the sheafification of the naive constant presheaf $A$ and identify the unit on global sections.
\end{exercise}

\begin{exercise}[Bounded continuous functions]
Show that the sheafification of the presheaf of bounded continuous real functions is the full sheaf of continuous real functions.
\end{exercise}

\begin{exercise}[A unit that is not injective]
For a space where every point has a proper neighborhood, analyze the presheaf $\mathcal F(X)=A$, $\mathcal F(U)=0$ for $U\subsetneq X$, and prove its sheafification is zero.
\end{exercise}

\begin{exercise}[Restriction to an open]
Let $j:U\hookrightarrow X$. Prove $(\mathcal F|_U)^\#\cong\mathcal F^\#|_U$.
\end{exercise}

\begin{exercise}[Cofinal basis neighborhoods]
Let $\mathcal B$ be a basis. Prove that the stalk $\mathcal F_x$ can be computed using only basis opens $B\ni x$.
\end{exercise}

\begin{exercise}[Separated presheaf]
Define separatedness for a presheaf and show every sheaf is separated.
\end{exercise}

\begin{exercise}[First plus is separated]
For the standard plus construction, prove that $\mathcal F^+$ is separated.
\end{exercise}

\begin{exercise}[Plus of a separated presheaf]
Assuming $\mathcal F$ is separated, prove that $\mathcal F^+$ satisfies gluing and is a sheaf.
\end{exercise}

\begin{exercise}[Two pluses]
Deduce that $\mathcal F^{++}$ is always a sheaf and has the universal property of sheafification.
\end{exercise}

\begin{exercise}[Adjunction]
Show that precomposition with the unit gives a natural bijection $\Hom_{\Sh(X)}(\mathcal F^\#,\mathcal G)\cong\Hom_{\PreSh(X)}(\mathcal F,\mathcal G)$ for every sheaf $\mathcal G$.
\end{exercise}

